\section{Respuestas a las Preguntas}

\textbf{1. Si la corriente de colector $I_C$ medida en la etapa BJT difiere significativamente del valor de diseño, ¿a qué parámetros del transistor se atribuye esta diferencia? Fundamentar considerando la dispersión de la ganancia de corriente ($\beta$ o $h_{FE}$) y de la tensión base-emisor ($V_{BE}$).}

La diferencia entre la corriente de colector ($I_C$) medida y el valor teórico se debe a que los cálculos analíticos asumen valores nominales. En la realidad física, parámetros intrínsecos como la ganancia de corriente ($\beta$ o $h_{FE}$) y la tensión base-emisor ($V_{BE}$) presentan una notable dispersión originada por tolerancias de fabricación y efectos térmicos.

Por un lado, la ganancia $\beta$ es el parámetro más inestable del dispositivo. En topologías de polarización dependientes, la corriente está dada por $I_C = \beta \cdot I_B$, por lo que cualquier variación en $\beta$ impacta proporcionalmente en $I_C$. 

Por otro lado, la tensión $V_{BE}$ teórica de $0.7\text{ V}$ varía en la práctica entre $0.6\text{ V}$ y $0.75\text{ V}$, experimentando además una deriva térmica de aproximadamente $-2\text{ mV/}^{\circ}\text{C}$. En una red polarizada por divisor de tensión, donde $I_C \approx \frac{V_{BB} - V_{BE}}{R_E}$, cualquier fluctuación en $V_{BE}$ produce un cambio porcentual severo en $I_C$, especialmente si el voltaje de Thévenin en la base ($V_{BB}$) es de baja magnitud.


\textbf{2. En la etapa de entrada formada por $Q_1$ y $Q_2$, ¿cómo influye la configuración Darlington en la estabilidad del punto de trabajo DC ($Q$) ante variaciones de temperatura y dispersión de $\beta$?}

La configuración Darlington ($Q_1$ y $Q_2$) compromete la estabilidad del punto de trabajo en continua ($Q$) debido a su alta sensibilidad térmica y a la severa multiplicación de tolerancias inherentes a los semiconductores.

En primer lugar, frente a las variaciones de temperatura, el arreglo coloca dos uniones base-emisor en serie ($V_{BE(total)} = V_{BE1} + V_{BE2}$). Puesto que cada unión presenta un coeficiente térmico de aproximadamente $-2\text{ mV/}^{\circ}\text{C}$, el conjunto duplica esta deriva a $-4\text{ mV/}^{\circ}\text{C}$, haciendo que el punto $Q$ sea inherentemente más inestable que en una etapa de transistor simple. Adicionalmente, la corriente de fuga térmica del primer transistor ($I_{CBO1}$) ingresa directamente a la base del segundo dispositivo ($Q_2$), el cual la multiplica por su propia ganancia ($\beta_2$). Este fenómeno provoca que, ante incrementos térmicos, la corriente de colector en reposo se desplace significativamente, pudiendo inducir un embalamiento térmico en casos extremos.

Por otro lado, respecto a la dispersión de la ganancia de corriente, el valor total del arreglo resulta del producto de las ganancias individuales ($\beta_{total} \approx \beta_1 \cdot \beta_2$). Esto provoca que las tolerancias de fabricación se multipliquen, generando una variabilidad masiva del parámetro $\beta_{total}$ entre distintos componentes. Debido a esta ganancia combinada tan elevada, la corriente de base ($I_B$) requerida resulta minúscula; en consecuencia, si la red no posee una topología de polarización solida e independiente de $\beta$, el punto $Q$ quedará totalmente a vulnerable a la dispersión del $\beta$.


\textbf{3. ¿Qué función cumple la resistencia $R_{Ein}$ en el emisor de $Q_2$? ¿Qué sucederá con el punto de polarización DC y la excursión de señal si $R_{Ein} = 0$?}

La resistencia $R_{Ein}$ opera como la carga estática de la etapa Darlington. Su función principal radica en establecer el punto de reposo en corriente continua (DC), limitando y fijando la magnitud de la corriente de colector ($I_{CQ}$) al proporcionar un camino resistivo hacia tierra.

En el escenario donde dicha resistencia adopta un valor nulo, el emisor de $Q_2$ queda físicamente cortocircuitado a masa. Bajo esta condición, las uniones base-emisor de los transistores quedarán polarizadas directamente sin ninguna limitación en la malla de salida. En consecuencia, la corriente de colector tenderá a un valor extremo (limitado únicamente por la fuente de alimentación), lo que ocasionará la destrucción térmica de los dispositivos.

\textbf{4. ¿Por qué la ganancia de tensión de esta etapa BJT es sustancialmente mayor que la de un amplificador MOSFET equivalente operando con corrientes de polarización similares? Relacionar con las expresiones de $g_m$ de ambos dispositivos.}

La superioridad en la ganancia de tensión de una etapa BJT frente a un MOSFET equivalente radica en la eficiencia de su transconductancia ($g_m$). Dado que en configuraciones análogas la ganancia es directamente proporcional a este parámetro ($A_v \approx -g_m \cdot R_L$), la diferencia de comportamiento proviene del principio físico rector de cada dispositivo. Al operar bajo la misma corriente de polarización ($I_C = I_D = I_Q$), el BJT exhibe una característica de transferencia exponencial, cuya transconductancia se define como:
$$g_{m,BJT} = \frac{I_C}{V_T}$$
donde $V_T \approx 26\text{ mV}$ a temperatura ambiente. En contraste, el MOSFET obedece a una ley cuadrática regida por campo eléctrico, con una transconductancia expresada como:
$$g_{m,MOSFET} = \frac{2I_D}{V_{OV}}$$
siendo $V_{OV} = V_{GS} - V_{th}$ la tensión de sobreimpulso.

Relacionando matemáticamente ambas expresiones para una misma corriente de polarización $I_Q$, se obtiene el siguiente cociente:
$$\frac{g_{m,BJT}}{g_{m,MOSFET}} = \frac{V_{OV}}{2V_T}$$

Asumiendo un valor típico de operación de $V_{OV} = 0.3\text{ V}$, la relación resulta ser aproximadamente $5.8$. En consecuencia, el BJT ofrece la máxima transconductancia por unidad de corriente posible en semiconductores ($\approx 38.5\text{ mS/mA}$ frente a los típicos $3$ a $8\text{ mS/mA}$ del MOSFET). Esto se traduce en que, ante la misma variación de tensión de entrada, el transistor bipolar genera un cambio de corriente mucho más agresivo, logrando así una ganancia de tensión entre 5 y 10 veces mayor que su contraparte de efecto de campo.

\textbf{5. Analizar el efecto de carga de la segunda etapa ($Q_3$) sobre la primera ($Q_1 - Q_2$). ¿Qué impedancia de entrada presenta la etapa de $Q_3$ ($Z_{in2}$) y cómo afecta a la ganancia en vacío de la primera etapa?}

El análisis del efecto de carga de la segunda etapa ($Q_3$) sobre la primera ($Q_1 - Q_2$) requiere evaluar la impedancia de entrada de la etapa posterior ($Z_{in2}$). De acuerdo con los datos relevados, $Z_{in2}$ presenta un valor  medido de $6.7\text{ k}\Omega$. Por su parte, la primera etapa exhibe una impedancia de salida en vacío ($Z_{out1}$) sumamente baja, con valor de $5.89\ \Omega$ en medición.

Dado que $Z_{in2}$ es varios órdenes de magnitud superior a $Z_{out1}$, el efecto de carga ejercido por $Q_3$ resulta prácticamente despreciable. El impacto sobre la ganancia en vacío de la primera etapa ($A_{v1,0}$) se determina mediante la relación del divisor de tensión conformado por ambas impedancias. La ganancia efectivamente transferida (cargada), denominada $A_{v1}$, se expresa a través de la siguiente ecuación:
$$A_{v1} = A_{v1,0} \cdot \left( \frac{Z_{in2}}{Z_{out1} + Z_{in2}} \right)$$

Al sustituir en la expresión analítica los valores experimentales obtenidos ($Z_{out1} = 5.89\ \Omega$ y $Z_{in2} = 6700\ \Omega$), se determina el factor de transferencia de la red:
$$\frac{Z_{in2}}{Z_{out1} + Z_{in2}} = \frac{6700\ \Omega}{5.89\ \Omega + 6700\ \Omega} \approx 0.9991$$

En consecuencia, la primera etapa transfiere aproximadamente el $99.9\%$ de su señal de tensión a la entrada de $Q_3$. Se concluye así que la ganancia en vacío de la configuración Darlington apenas experimenta atenuación tras la conexión de la segunda etapa amplificadora.

\textbf{6. ¿Por qué modificar la resistencia de colector $R_C$ de la segunda etapa altera drásticamente la ganancia total del amplificador, mientras que modificar $R_{Ein}$ impacta principalmente sobre la impedancia de entrada?}

La alteración de la resistencia de colector ($R_C$) en la segunda etapa ($Q_3$) incide drásticamente sobre la ganancia total del amplificador debido a que dicha etapa opera en configuración de emisor común. En esta topología, la capacidad de amplificación de voltaje es directamente proporcional a la carga resistiva presente en el nodo del colector. Matemáticamente, la magnitud de la ganancia de tensión se rige por la siguiente relación de transferencia dinámica:
$$A_v \approx -g_m \cdot R_C$$
Al alterar $R_C$, se modifica directamente este factor de carga. Un incremento en dicha resistencia permite desarrollar una mayor caída de tensión alterna para una misma variación dinámica en la corriente de colector, alterando de manera directa y lineal la ganancia total de voltaje del sistema.

Por el contrario, la modificación de la resistencia $R_{Ein}$ en la primera etapa afecta primordialmente a la impedancia de entrada, dado que el par Darlington actúa como un seguidor de emisor cuya propiedad fundamental es la transformación de impedancias. La resistencia estática ubicada en el emisor ($R_{Ein}$) se refleja hacia la entrada del circuito (la base de $Q_1$) multiplicada por la ganancia de corriente de ambos semiconductores, aproximándose mediante la expresión:
$$Z_{in} \approx \beta_1 \cdot \beta_2 \cdot R_{Ein}$$
Al modificar $R_{Ein}$, el cambio cuantitativo es amplificado por el producto masivo de las ganancias de corriente ($\beta_1 \cdot \beta_2$). Por consiguiente, el impacto se manifiesta casi exclusivamente sobre la impedancia de entrada total de la red, lo cual determina el nivel de acoplamiento y la exigencia de corriente hacia la fuente de señal, sin aportar ninguna contribución a la ganancia de voltaje.

\textbf{7. Proponer dos modificaciones de diseño para aumentar la ganancia total $G_T$ sin alterar las corrientes de polarización. Evaluar su practicidad de implementación, su impacto en la tensión de excursión simétrica y en la disipación de potencia.}

Dado que las corrientes de polarización deben mantenerse inalteradas, la transconductancia ($g_m$) de los dispositivos permanece constante. Por lo tanto, el incremento de la ganancia total ($G_T$) se debe lograr manipulando exclusivamente la red dinámica de alterna (AC) en la segunda etapa ($Q_3$). La primera modificación consiste en el desacople de la resistencia de emisor ($R_E$) mediante la incorporación de un capacitor paralelo ($C_E$) de alta capacidad. En corriente continua (DC), el capacitor actúa como un circuito abierto, preservando el punto de reposo, la disipación de potencia y la tensión de excursión máxima simétrica sin alteraciones. En alterna, $C_E$ opera como un cortocircuito a masa, suprimiendo la realimentación negativa (degeneración de emisor) y elevando la ganancia a su valor intrínseco de lazo abierto:
$$A_v \approx -g_m \cdot R_C$$
Esta solución presenta una alta practicidad de implementación, siendo una maniobra estándar, económica y simple que requiere únicamente el agregado de un componente pasivo.

La segunda alternativa radica en reemplazar la resistencia de colector ($R_C$) por una carga activa, implementada convencionalmente mediante un espejo de corriente con transistores PNP. Esta red inyecta la corriente de polarización nominal exacta, manteniendo el punto de reposo DC inalterado, pero exhibe una impedancia dinámica ($r_o$) extremadamente elevada frente a señales de pequeña amplitud. En consecuencia, la ganancia de tensión se dispara a niveles máximos según la expresión:
$$A_v \approx -g_m \cdot r_o$$
Sin embargo, su practicidad es moderada a baja en ensambles con componentes discretos, dado que exige semiconductores adicionales rigurosamente apareados. En cuanto a sus impactos, si bien la disipación de potencia estática global no varía de forma significativa, la excursión simétrica se ve reducida; el transistor de la carga activa requiere una caída de tensión mínima (voltaje de saturación, $V_{CE(sat)}$) para operar en su región lineal, lo que sustrae margen de voltaje respecto al riel de alimentación y comprime el límite antes de alcanzar el recorte de la señal.

\textbf{8. La especificación requiere $Z_{in} > 350\text{ k}\Omega$. ¿Por qué resulta indispensable utilizar el par Darlington ($Q_1 - Q_2$) en lugar de una etapa Emisor Común convencional para cumplir este requerimiento?}

Para alcanzar una impedancia de entrada superior a $350\text{ k}\Omega$, resulta inviable emplear una etapa de emisor común con un único transistor debido a sus limitaciones estructurales. En dicha topología, la impedancia base depende del parámetro $h_{ie}$ y de la reflexión de la resistencia de emisor, aproximándose como $Z_{in} \approx \beta \cdot R_E$. Dado que la ganancia de corriente ($\beta$) estándar oscila entre 100 y 300, satisfacer la especificación exigiría un valor de $R_E$ tan elevado que demandaría voltajes de alimentación impracticables o corrientes de polarización microscópicas, degradando severamente la linealidad del sistema; esto sin mencionar que el divisor resistivo de polarización en paralelo mermaría aún más la impedancia neta. Por el contrario, el uso del par Darlington configurado como colector común (seguidor de emisor) se torna indispensable porque multiplica las ganancias de corriente de ambos semiconductores ($\beta_{total} \approx \beta_1 \cdot \beta_2$). Esta multiplicación extrema permite que la carga dinámica en el emisor se refleje hacia la entrada por un factor masivo, elevando la impedancia intrínseca del circuito al rango de los mega-ohmios mediante la relación:
$$Z_{in} \approx \beta_1 \cdot \beta_2 \cdot R_E$$
De este modo, el acoplamiento supera ampliamente el requerimiento de diseño sin comprometer la estabilidad del punto de reposo ni las características de transferencia de la etapa amplificadora.

\textbf{9. Si el generador de funciones del laboratorio presenta una impedancia de salida de $50\ \Omega$, ¿por qué la caída de tensión en dicha resistencia interna es despreciable frente a la $Z_{in}$ medida?}

La interacción entre la fuente de señal y el amplificador se rige por el principio del divisor de tensión, donde la resistencia interna del generador ($R_g = 50\ \Omega$) y la impedancia de entrada del amplificador ($Z_{in}$) quedan configuradas en serie. La tensión efectiva que ingresa a los terminales de la primera etapa ($V_{in}$) se determina a partir de la siguiente ecuación de transferencia:
$$V_{in} = V_g \cdot \left( \frac{Z_{in}}{R_g + Z_{in}} \right)$$
donde $V_g$ representa la tensión electromotriz ideal de la fuente. Dado que la impedancia de entrada lograda por la topología Darlington (del orden de los cientos de kilo-ohmios) es masivamente superior a la resistencia del generador ($Z_{in} \gg R_g$), el denominador de la expresión está dominado enteramente por $Z_{in}$. Al evaluar el divisor, el factor de acoplamiento resulta prácticamente unitario ($\approx 0.9998$). Matemáticamente, esto implica que más del $99.98\%$ de la señal se transfiere de forma íntegra hacia la base del primer transistor, limitando la caída de tensión sobre la resistencia interna a un nivel inferior al $0.02\%$. En consecuencia, esta pérdida resulta de una magnitud absolutamente despreciable, garantizando que la etapa de entrada opere de forma análoga a un voltímetro ideal, es decir, sensando la tensión sin introducir atenuación ni demandar corriente significativa al instrumental de laboratorio.

\textbf{10. La especificación exige una impedancia de salida $Z_o = 8\ \Omega \pm 5\%$ para alimentar una carga $R_L = 8\ \Omega$. ¿Qué característica del circuito permite lograr un valor tan bajo de $Z_o$ y qué ocurriría con la ganancia si se conectara $R_L$ directamente al colector de un Emisor Común?}

El cumplimiento de esta especificación se sustenta en el diseño de la etapa de potencia, la cual emplea un par Sziklai configurado funcionalmente como un amplificador de corriente. Esta topología permite dividir la impedancia proveniente de la etapa amplificadora de tensión previa por el producto de las ganancias de corriente de los transistores de potencia ($\beta_1 \cdot \beta_2$). Como resultado, se obtiene una impedancia de salida intrínseca del semiconductor extraordinariamente baja, del orden de los miliohmios, lo cual permite al diseñador fijar con precisión la impedancia de salida final del amplificador ($Z_o = 8\ \Omega$) mediante la adición de componentes resistivos pasivos en serie con la carga. Por el contrario, si la carga de baja impedancia ($R_L = 8\ \Omega$) se conectara directamente al nodo del colector de una etapa en emisor común, la ganancia de tensión del sistema colapsaría. En dicha topología, la ganancia es directamente proporcional a la resistencia de carga efectiva en alterna, rigiéndose por la ecuación:
$$A_v \approx -g_m \cdot (R_C \parallel R_L)$$
Al conectar una carga tan exigente, el paralelo equivalente se desploma a un valor cercano a los $8\ \Omega$ (independientemente del valor de $R_C$), destruyendo la capacidad del circuito para desarrollar una amplitud de voltaje significativa y atenuando severamente la señal.

\textbf{11. Al incrementar paulatinamente el nivel de tensión de entrada $v_{in}$, la salida comienza a distorsionarse. ¿Qué mecanismo del BJT limita la excursión máxima por el extremo superior (positivo) y cuál por el extremo inferior (negativo) en la etapa de salida?}

La distorsión asimétrica en los extremos de la señal se explica por la naturaleza intrínseca del transistor bipolar, cuyo comportamiento de transferencia emula la curva exponencial de un diodo ($i_C \approx I_S \cdot e^{v_{BE}/V_T}$). Por el extremo que conduce a la saturación, la excursión está limitada por el incremento drástico de la conducción; al elevarse la excitación, el punto de operación alcanza una zona donde la pendiente de la curva exponencial de la unión base-emisor se vuelve extremadamente pronunciada (casi vertical). En esta región, un incremento infinitesimal en la tensión de entrada exige un aumento masivo de corriente que la red de polarización no puede suministrar, forzando a que la tensión colector-emisor colapse abruptamente a su límite físico mínimo ($V_{CE(sat)}$) y achatando el pico de la onda. Por el extremo opuesto, la limitación es gobernada por el mecanismo de corte; cuando el semiciclo de la señal reduce la polarización por debajo del umbral de conducción de la unión (ingresando a la región donde la curva exponencial del diodo se aplana asintóticamente a cero), la corriente de colector se anula ($I_C = 0$). Bajo esta condición, el dispositivo deja de conducir y la tensión de salida queda rígidamente enclavada, cercenando el otro extremo de la señal alterna.



\textbf{12. A partir de la relación exponencial $i_C = I_S \cdot e^{v_{be}/V_T}$, explicar por qué el BJT genera distorsión armónica ante amplitudes de entrada pequeñas (decenas de mV) y cómo contribuye la resistencia de emisor no puenteada ($R_{E1}$) a linealizar la respuesta.}

La generación de distorsión armónica en un transistor BJT tiene su origen en la naturaleza no lineal de su característica de transferencia, regida por la relación exponencial $i_C = I_S \cdot e^{v_{be}/V_T}$. Al someter la unión a una señal senoidal y expandir la función mediante series de Taylor, la corriente resultante se compone de un término fundamental y de múltiples sumandos de orden superior que originan los armónicos. Para que dichas componentes sean despreciables, la amplitud debe respetar estrictamente la condición de pequeña señal ($v_{be} \ll V_T \approx 26\text{ mV}$); sin embargo, cuando la excitación alcanza decenas de milivoltios, el cociente $\frac{v_{be}}{V_T}$ adquiere un peso matemático significativo, deformando la onda de salida. Para mitigar este fenómeno, la inclusión de una resistencia de emisor no puenteada ($R_{E1}$) introduce una realimentación negativa local que modifica la transconductancia efectiva de la etapa según la expresión:
$$g_m' = \frac{g_m}{1 + g_m R_{E1}}$$
Al garantizar por diseño que $g_m R_{E1} \gg 1$, este parámetro dinámico se independiza de la no linealidad del semiconductor y se aproxima a $g_m' \approx \frac{1}{R_{E1}}$. En consecuencia, la transferencia de la etapa ($i_c \approx \frac{v_{in}}{R_{E1}}$) pasa a estar gobernada exclusivamente por un componente pasivo lineal, lo que suprime drásticamente la distorsión armónica del sistema.

\textbf{13. ¿A partir de qué amplitud de tensión de entrada la distorsión de la señal se hizo perceptible en el osciloscopio? Relacionar dicha medición con el valor de THD obtenido y el umbral de detección visual ($3\% \text{--} 5\%$).}

Durante el ensayo de laboratorio, la deformación de la señal se hizo perceptible a simple vista al alcanzar una tension de $18\text{ mV}_{pp}$. Dado que el protocolo de ensayo fue estipulado el registro de la Distorsión Armónica Total (THD) en niveles de tensión predefinidos, la caracterización espectral se tomó sobre el escalón de los $20\text{ mV}_{pp}$, arrojando un valor de distorsión del $13\%$. Esta correlación concuerda con las limitaciones de resolución del ojo humano. Por consiguiente, se infiere que al cruzar la frontera de los $18\text{ mV}_{pp}$, el contenido armónico de la señal rompió dicha barrera crítica del $5\%$, volviéndose evidente.

\textbf{14. ¿Qué componente del circuito presenta la mayor disipación de potencia en reposo ($P_{DC}$)? Justificar mediante las mediciones de $V_{CE}$ e $I_C$, y verificar si opera dentro de su Zona de Operación Segura (SOA).}

El análisis del balance energético indica que, si bien la resistencia $R_{ES}$ registra el mayor consumo absoluto del circuito con $6.72\text{ W}$[cite: 1], el componente activo que presenta la mayor disipación de potencia en reposo ($P_{DC}$) es el transistor $Q_5$, alcanzando un valor medido de $3.08\text{ W}$. Esta magnitud térmica se justifica analíticamente mediante el producto de sus coordenadas estáticas de polarización, calculándose como $P_{DC} = V_{CE} \cdot I_C$. Al proyectar estos parámetros operativos (correspondientes a un $V_{CE}$ de $3.95\text{ V}$ y una $I_C$ de $0.779\text{ A}$) sobre la gráfica de la Zona de Operación Segura (SOA) estipulada en la hoja de datos del fabricante, se constata que el punto de trabajo estático se ubica firmemente dentro de los márgenes térmicos y eléctricos permitidos. Esto garantiza que el dispositivo opera exento de riesgos por ruptura o fatiga térmica, dado que cuenta con un disipador térmico  adecuado para evacuar los $3.08\text{ W}$ generados hacia el ambiente.

\textbf{15. Al reducir la tensión de alimentación $V_{CC}$ por debajo de los $12\text{ V}$ nominales, ¿cuál es el primer parámetro que se degrada (ganancia de tensión, excursión máxima sin recorte o impedancias)? Fundamentar.}

Al reducir la tensión de alimentación ($V_{CC}$) por debajo de su valor nominal, el primer parámetro que experimenta una degradación drástica es la excursión máxima de señal. Este fenómeno se debe a que la amplitud máxima de la onda está físicamente delimitada por el nivel de tensión de los rieles de la fuente; al disminuir $V_{CC}$, el espacio disponible para que la señal oscile se comprime, provocando que la onda alcance prematuramente las zonas de saturación o corte y sufra un recorte severo. 

